% !TEX program = xelatex
\documentclass[12pt, a4paper]{article}
\usepackage{fontspec}
\usepackage{polyglossia}
\setdefaultlanguage[locale=mashriq]{arabic}
\setotherlanguage{english}
\newfontfamily\arabicfont[Script=Arabic, Scale=1.1]{Amiri}
\newfontfamily\englishfont[Script=Latin]{TeX Gyre Pagella}

\usepackage[backend=biber, style=authoryear, sorting=nyt]{biblatex}
\addbibresource{references.bib}

\begin{document}

\section{مراجع عربية في BibLaTeX}

يُوضّح هذا المستند كيفية إدراج مراجع عربية في قائمة مراجع
BibLaTeX مع دعم UTF-8 الكامل.

كما أشار \textcite{alKhwarizmi1983}، وُضعت أسس الجبر في القرن
التاسع الميلادي. وفي وقت لاحق، أسهم \textcite{ibnSina2005} في
العلوم الطبية.

\printbibliography

\end{document}
